\documentclass[11pt,a4paper]{article}

% ------------------------------------------------------------
% 한글: XeLaTeX로 컴파일 (Overleaf에서는 Compiler를 XeLaTeX로 설정)
% 로컬에서 kotex가 있으면 fontspec 대신 \usepackage{kotex}를 써도 된다.
% ------------------------------------------------------------
\usepackage{fontspec}
\setmainfont{Noto Serif CJK KR}
\setsansfont{Noto Sans CJK KR}
\setmonofont{DejaVu Sans Mono}[Scale=0.85]
\XeTeXlinebreaklocale "ko"
\XeTeXlinebreakskip = 0pt plus 1pt

\usepackage[margin=2.3cm]{geometry}
\usepackage{amsmath,amssymb}
\usepackage{booktabs,tabularx,array,multirow,longtable}
\usepackage{enumitem}
\usepackage{xcolor}
\usepackage{float}
\usepackage{tikz}
\usetikzlibrary{positioning,arrows.meta}
\usepackage[hidelinks]{hyperref}

\setlength{\parskip}{0.5em}
\setlength{\parindent}{0pt}
\renewcommand{\arraystretch}{1.25}
\newcolumntype{Y}{>{\raggedright\arraybackslash}X}

\definecolor{boxbg}{RGB}{243,242,238}
\newcommand{\keypoint}[1]{%
  \par\smallskip\noindent\colorbox{boxbg}{\parbox{\dimexpr\linewidth-2\fboxsep}{#1}}\par\smallskip}

\title{5PL 미들마일 플랫폼의 동적 매칭·잉여배분 연구\\[4pt]
\large 모델, 알고리즘, 실험 설계와 현재 진행 상황}
\author{APIEMS 2026 연구 진행 문서}
\date{2026년 9월 22일 기준}

\renewcommand{\contentsname}{차례}
\renewcommand{\figurename}{그림}
\renewcommand{\tablename}{표}
\renewcommand{\appendixname}{부록}
\begin{document}
\maketitle

\keypoint{\textbf{한 줄 요약.}
주문자·공급자를 매칭하고 거래잉여를 나누는 플랫폼에서, 참여자를 붙잡는 가치(유지 가치)를
학습해 매 기간 최적화에 반영하는 방법을 연구한다.
현재 슬라이스 0--4(시장 구축, 비교군, 선형 근사, 신경망 근사와 실험 1)를 마쳤고,
다음은 그래프 신경망을 이용한 실험 3이다.
실험 1에서는 가설과 달리 신경망이 선형 근사보다 낫지 않았고, 두 근사 정책 모두 근시안 정책을 넘지 못했다.}

\tableofcontents
\newpage

% ============================================================
\section{연구 질문}
% ============================================================

플랫폼의 오늘 결정(어떤 주문을 받고, 누구에게 맡기고, 이익을 어떻게 나누는지)은
내일 누가 플랫폼에 남아 있는지를 바꾼다. 따라서 좋은 결정을 하려면
``이 참여자를 붙잡아 두면 앞으로 얼마나 이득인가''를 알아야 한다.

\keypoint{\textbf{연구 질문.} 현재의 매칭·잉여배분이 미래 참여자 구성을 바꾸는 플랫폼에서,
미래 가치를 어떻게 근사해야 동적 의사결정이 개선되는가?}

미래 가치를 정확히 계산하는 것은 불가능하므로 근사해야 한다. 기존 방식(선형 근사)은
참여자 수, 품목별 공급량 같은 요약 수치를 따로따로 더하기 때문에 두 가지를 놓친다.
첫째, 수치들이 \emph{함께} 작용하는 효과(예: 묶음 주문은 모든 품목이 있어야 성립)를 표현하지 못한다.
둘째, 요약 수치가 같아도 \emph{누가 어떤 품목을 함께 공급하는지}에 따라
특정 공급자가 떠났을 때의 피해가 달라진다는 점을 구분하지 못한다.
이 연구는 신경망과 그래프 신경망으로 이 두 한계를 완화할 수 있는지 검증한다.

% ============================================================
\section{문제 상황과 모델}
% ============================================================

\subsection{플랫폼과 참여자}

\begin{itemize}[leftmargin=1.5em]
  \item \textbf{주문자}: 생산용 부품을 여러 품목의 \emph{묶음}으로 주문한다
        (예: 품목 A 10개 + 품목 B 8개). 품목별 수량과 제안가격을 가진다.
  \item \textbf{공급자}: 몇 개의 품목을 공급하며, 품목별 공급가능량(용량)과 공급가격을 가진다.
  \item \textbf{플랫폼}: 매 기간 주문을 받을지, 품목별로 어느 공급자에게 몇 개씩 맡길지,
        거래로 생긴 잉여를 주문자·공급자·플랫폼이 얼마씩 나눌지 정한다.
        물건은 단일 허브에 모아 주문자에게 한 번에 배송한다.
\end{itemize}

두 가지 규칙이 문제의 핵심이다.
\begin{enumerate}[leftmargin=1.5em]
  \item \textbf{묶음 성립 규칙}: 주문은 요청한 \emph{모든} 품목의 요구수량이 확보될 때만 성립한다.
        한 품목이라도 부족하면 주문 전체가 무산된다.
  \item \textbf{재참여 규칙}: 참여자가 다음 기간에 다시 참여할 확률은
        이번 기간에 배분받은 잉여에 따라 달라진다.
\end{enumerate}

\subsection{모델 가정}

\begin{itemize}[leftmargin=1.5em]
  \item \textbf{참여자 프로필}: 주문자는 고유한 품목 구성·수량·가격 수준을,
        공급자는 고유한 공급 품목·공급량·가격 수준을 가진다.
        재참여하면 프로필에 작은 변동을 더해 주문·공급한다.
  \item \textbf{분할 공급 허용}: 한 주문의 한 품목을 여러 공급자가 나눠 공급할 수 있다
        (허브에서 모아 통합 배송하므로).
  \item \textbf{배송비}: 성립한 주문마다 허브$\to$주문자 배송 고정비 $f=25$를 한 번 부과한다.
        수량에 따른 운송·작업비는 제안가격과 공급가격의 차이(잉여)에 이미 반영된 것으로 본다.
  \item \textbf{기회손실}: 성립하지 못한 주문의 놓친 이윤은 지표로만 기록하고 이윤에서 빼지 않는다.
        거절의 미래 손해는 유지 가치로 반영되므로 이중 계산을 피한다.
  \item \textbf{플랫폼 이윤 제약}: 기간 전체 합 기준으로 음수가 되지 않게 한다.
        같은 기간 안의 교차 보조(유지 가치가 큰 참여자를 위해 한 거래에서 손해를 감수하고
        다른 거래 이윤으로 메우는 것)는 허용한다. 기간을 넘는 적자는 허용하지 않는다.
  \item \textbf{주문자 유입}: 매 기간 신규 주문자가 평균 10명 도착한다.
  \item \textbf{공급자 자리}: 공급자 자리는 6개로 고정되어 있다. 공급자가 떠나면 그 자리는
        매 기간 확률 $0.25$로(평균 4기간 뒤) \emph{같은 품목 조합}의 새 공급자로 채워진다.
        새 공급자를 플랫폼에 들이는 데(검증·계약·등록) 시간이 걸리는 현실을 반영하며,
        충원이 즉시 이루어지면 공급자를 잃는 손해가 사라져 유지 가치가 의미를 잃는다.
  \item \textbf{거래 이력 제외}: 재참여 확률은 해당 기간의 배분 잉여에만 의존하며,
        누적 거래 이력(충성도 효과)은 이번 범위에 포함하지 않는다.
\end{itemize}

\subsection{시장 조건}

실험 3에서 바꾸는 세 가지 값을 \textbf{시장 조건}이라 부른다.

\begin{itemize}[leftmargin=1.5em]
  \item \textbf{공급 편중도} ($\mathrm{conc}$): 공급자들이 품목을 얼마나 고르게 나눠 갖는지.
        공급자 수, 품목 수, 품목당 대체 공급자 수를 모두 고정하고, 공급자별 품목 수의 분포만 바꾼다.
        0이면 모두 같은 수의 품목을 맡고, 1에 가까울수록 소수 공급자에게 품목이 몰린다.
  \item \textbf{대체 공급자 수} ($n_{\mathrm{alt}}$): 한 품목을 공급할 수 있는 공급자 수.
  \item \textbf{주문당 품목 수} ($k$): 주문 하나에 묶인 품목 수.
\end{itemize}

공급자 수를 고정했으므로 공급자당 평균 품목 수는
$(\text{품목 수}\times n_{\mathrm{alt}})/(\text{공급자 수})$로 정해진다.
기본 규모(품목 6, 공급자 6, $n_{\mathrm{alt}}=2$)에서 평균은 2이다.

\begin{table}[H]
\centering\small
\caption{공급 편중도별 공급자별 품목 수 (실제 생성 결과)}
\begin{tabular}{cccccc}
\toprule
$n_{\mathrm{alt}}$ & 평균 품목 수 & $\mathrm{conc}=0$ & $\mathrm{conc}=0.5$ & $\mathrm{conc}=1$ & 비고\\
\midrule
1 & 1 & [1,1,1,1,1,1] & 같음 & 같음 & 편중도 변경 불가\\
2 & 2 & [2,2,2,2,2,2] & [4,2,2,2,1,1] & [6,2,1,1,1,1] & 핵심 설계\\
3 & 3 & [3,3,3,3,3,3] & [5,4,3,2,2,2] & [6,6,3,1,1,1] & 보조 설계\\
\bottomrule
\end{tabular}
\end{table}

\keypoint{\textbf{핵심 성질.} 공급 편중도를 바꿔도 공급자 수, 품목별 대체 공급자 수,
품목별 총 공급용량, 기대 수요, 그리고 시장 요약 지표가 모두 같다. 달라지는 것은
``어떤 공급자가 몇 개의 품목을 떠받치는가''라는 구조뿐이다.}

주문당 품목 수 $k$를 바꿀 때는 주문자 수를 고정하고 품목당 주문 수량을 $k_{\mathrm{ref}}/k$배
($k_{\mathrm{ref}}=2$)로 조정해, 주문 한 건의 총량과 품목별 기대 수요를 같게 유지한다.

\begin{table}[H]
\centering\small
\caption{주문당 품목 수에 따른 주문 수량 (수량 단위 3배 적용)}
\begin{tabular}{cccccc}
\toprule
$k$ & 품목당 주문 수량 & 주문당 평균 총량 & 품목별 기대 수요 & 공급 여유 & 공급자 구조\\
\midrule
1 & 12--30 & 21 & 같음 & 같음 & 같음\\
2 & 6--15 & 21 & 같음 & 같음 & 같음\\
3 & 4--10 & 21 & 같음 & 같음 & 같음\\
\bottomrule
\end{tabular}
\end{table}

\subsection{재참여 모형}

참여자 $n$이 배분받은 잉여를 $s_n$, 그 참여자의 \textbf{제안 금액}을 $w_n$이라 할 때,
재참여 확률은 잉여율 $s_n/w_n$의 로지스틱 함수다.
\begin{equation}
  P(\text{재참여}) = \sigma\!\left(a + b_{\mathrm{kind}}\cdot\frac{s_n}{w_n}\right),
  \qquad \sigma(z)=\frac{1}{1+e^{-z}}
\end{equation}
제안 금액은 주문자와 공급자에 대해 다음과 같다.
\begin{align}
  w_b &= \textstyle\sum_i p_{bi}\,q_{bi} && \text{(주문자 $b$: 주문 전체 금액)}\\
  w_j &= \textstyle\sum_i c_{ji}\,\mathrm{cap}_{ji} && \text{(공급자 $j$: 내놓은 공급 전체 금액)}
\end{align}
공급자의 분모는 실제로 판 양이 아니라 \emph{내놓은 공급 전체}이므로,
여러 품목을 많이 내놓은 공급자는 같은 잉여를 받아도 잉여율이 낮아 떠나기 쉽다.
이는 논문에서 명시해야 할 모델링 선택이다.

곡선은 뜻이 분명한 두 점으로 직접 정했다(주문자·공급자 동일).
\begin{itemize}[leftmargin=1.5em]
  \item 잉여 0일 때 재참여 확률 $0.5$ (거래하지 못해도 절반은 남는다) $\Rightarrow a=0$
  \item 기준 잉여율(기준 몫 $(0.3,0.3)$으로 운영할 때의 평균 잉여율)에서 재참여 확률 $0.85$
\end{itemize}
기준 잉여율은 튜닝 seed에서 한 번 실측했으며(주문자 $0.1048$, 공급자 $0.1020$),
이로부터 기울기는 $b_{\mathrm{buy}}=16.55$, $b_{\mathrm{sup}}=17.01$이다.

\subsection{공통 난수}

정책끼리 공정하게 비교하려면 모든 정책이 \emph{같은 운}을 받아야 한다.
난수는 (seed, 반복, 스트림, 참여자, 기간)을 키로 뽑는다.
재참여는 참여자·기간별 균등난수 $u$를 미리 뽑아 두고 $u<P(\text{재참여})$로 판정한다.
따라서 정책이 달라도 같은 참여자·같은 기간에는 같은 난수가 쓰이고,
신규 도착과 변동은 이전 결정과 무관하게 같다.

\subsection{규모와 주요 파라미터}

\begin{table}[H]
\centering\small
\caption{현재 설정 (\texttt{utils/params.py}의 \texttt{Cfg})}
\begin{tabularx}{\linewidth}{lY}
\toprule
항목 & 값\\
\midrule
품목, 공급자 자리, 대체 공급자 수 & 6, 6, 2 (공급자 수는 시장 조건과 무관하게 고정)\\
기본 시장 조건 & $\mathrm{conc}=0.5$, $k=2$\\
신규 주문자 & 기간당 평균 10명\\
품목별 총 공급용량 & $1.3\times$ 안정 상태 수요 $\div$ 실측 공급자 점유율 $0.602$ $\approx 151$\\
공급자 빈자리 충원 확률 & 0.25 (평균 4기간)\\
배송 고정비 & 25 (주문당)\\
기간 $T$, 평가 반복 수 & 30, 30\\
재참여 곡선 & $P(0)=0.5$, $P(\text{기준 잉여율})=0.85$\\
규칙 기반 몫 & 주문자 0.3, 공급자 0.3 (격자 튜닝 결과 기준값 유지)\\
구간선형 근사 & 구간 5개, 최대 오차 $\le 0.03$\\
솔버 & Gurobi 13.0.3 (학술 라이선스), Seed·Threads·MIPGap 고정\\
seed 분리 & 튜닝·학습 seed 1, 평가 seed 0\\
\bottomrule
\end{tabularx}
\end{table}

% ============================================================
\section{한 기간의 흐름}
% ============================================================

\begin{figure}[H]
\centering
\begin{tikzpicture}[
  box/.style={draw, rounded corners=3pt, minimum height=0.85cm, align=center, font=\small},
  arr/.style={-{Stealth[length=2.2mm]}, thick}]
\node[box, minimum width=6.4cm] (s) at (0,0) {현재 상태: 활동 중인 주문자·공급자와 프로필};
\node[box, minimum width=4.0cm] (f) at (-3.2,-1.5) {시장 요약 지표\\ \footnotesize 선형 · MLP · MLP+편중도 요약};
\node[box, minimum width=4.0cm] (g) at (3.2,-1.5) {거래 연결 구조 (그래프)\\ \footnotesize GNN};
\node[box, minimum width=6.4cm] (v) at (0,-3.0) {미래 가치 예측 $\hat V(s)$};
\node[box, minimum width=6.4cm] (c) at (0,-4.3) {참여자별 유지 가치 $c_n = \hat V(s)-\hat V(s\setminus n)$};
\node[box, minimum width=6.4cm] (m) at (0,-5.6) {MILP: 수락 · 배정 · 잉여배분};
\node[box, minimum width=6.4cm] (e) at (0,-6.9) {환경: 거래 실행 (묶음 성립 검사)};
\node[box, minimum width=6.4cm] (r) at (0,-8.2) {재참여 추첨 (로지스틱, 공통 난수)};
\draw[arr] (s) -- (f); \draw[arr] (s) -- (g);
\draw[arr] (f) -- (v); \draw[arr] (g) -- (v);
\draw[arr] (v) -- (c); \draw[arr] (c) -- (m); \draw[arr] (m) -- (e); \draw[arr] (e) -- (r);
\draw[arr] (r.west) -- (-6.0,-8.2) -- node[pos=0.5, left, font=\footnotesize, align=right]{다음\\기간} (-6.0,0) -- (s.west);
\end{tikzpicture}
\caption{한 기간의 결정 흐름. 근시안 정책은 유지 가치를 0으로, 규칙 기반 정책은 잉여를 고정 비율로 둔다.}
\end{figure}

모든 정책은 \textbf{같은 MILP}를 쓰고, 유지 가치 $c_n$을 어디서 받아오는지만 다르다.
따라서 정책 간 차이는 오직 미래 가치를 얼마나 잘 예측했는지에서 나온다.

% ============================================================
\section{기간별 최적화 모형 (MILP)}
% ============================================================

아래 식은 \texttt{src/match/milp\_build.py}의 구현을 그대로 옮긴 것이다.

\subsection{집합과 파라미터}

\begin{tabularx}{\linewidth}{lY}
$b \in \mathcal{B}$ & 이번 기간 활동 중인 주문(주문자)\\
$j \in \mathcal{J}$ & 이번 기간 활동 중인 공급자\\
$i \in \mathcal{I}$ & 품목\\
$q_{bi}$ & 주문 $b$의 품목 $i$ 요구수량\\
$\mathrm{cap}_{ji}$ & 공급자 $j$의 품목 $i$ 공급가능량\\
$p_{bi}$ & 주문 $b$의 품목 $i$ 제안가격 \\
$c_{ji}$ & 공급자 $j$의 품목 $i$ 공급가격\\
$f$ & 주문당 배송 고정비 ($=25$)\\
$w_b, w_j$ & 주문자·공급자의 제안 금액 (재참여 잉여율의 분모)\\
$c_b, c_j$ & 주문자·공급자의 유지 가치 (근사 모델이 제공, 근시안은 0)\\
$(\alpha,\beta)$ & 규칙 기반 정책의 주문자·공급자 몫\\
\end{tabularx}

\subsection{결정변수}

\begin{align*}
  x_{bji} &\in \mathbb{Z}_{\ge 0},\ \ x_{bji}\le \min(q_{bi},\mathrm{cap}_{ji}) && \text{공급자 $j$가 주문 $b$의 품목 $i$에 보내는 수량}\\
  y_b &\in \{0,1\} && \text{주문 $b$ 성립 여부}\\
  u_b &\ge 0 && \text{주문자 $b$에게 배분하는 잉여 (주문 단위)}\\
  v_j &\ge 0 && \text{공급자 $j$에게 배분하는 잉여 (공급자 단위)}
\end{align*}

마진을 다음과 같이 정의한다.
\begin{align}
  g_b &= \sum_{j}\sum_{i}(p_{bi}-c_{ji})\,x_{bji} && \text{(주문 $b$의 거래 마진)}\\
  m_j &= \sum_{b}\sum_{i}(p_{bi}-c_{ji})\,x_{bji} && \text{(공급자 $j$가 관여한 거래 마진)}
\end{align}

\subsection{목적함수}

\begin{equation}
\begin{aligned}
  \max\ & \underbrace{\sum_b g_b - f\sum_b y_b - \sum_b u_b - \sum_j v_j}_{\text{플랫폼 당기 이윤 } \pi}\\
  &+\; \underbrace{\sum_b c_b\,\tilde P_{\mathrm{buy}}(u_b; w_b) + \sum_j c_j\,\tilde P_{\mathrm{sup}}(v_j; w_j)}_{\text{유지 가치 항}}
\end{aligned}
\end{equation}

$\tilde P(\cdot\,;w)$는 재참여 확률의 구간선형 근사로, 잉여율 꺾인 점에 $w$를 곱해
참여자마다 꺾인 점을 만들고 Gurobi의 \texttt{addGenConstrPWL}로 표현한다.
유지 가치가 높은 참여자일수록 잉여를 더 줘서 재참여 확률을 올리는 것이 이득이 된다.

\subsection{제약식}

\begin{align}
  &\sum_{j} x_{bji} = q_{bi}\,y_b && \forall b,i && \text{(묶음 성립·분할 공급)}\label{c:bundle}\\
  &\sum_{b} x_{bji} \le \mathrm{cap}_{ji} && \forall j,i && \text{(공급자 용량)}\label{c:cap}\\
  &u_b = \alpha\, g_b,\quad v_j = \beta\, m_j && && \text{(규칙 기반 정책)}\label{c:rule}\\
  &u_b \le g_b,\quad v_j \le m_j && && \text{(그 외 정책)}\label{c:free}\\
  &\sum_b g_b - f\sum_b y_b - \sum_b u_b - \sum_j v_j \ge 0 && && \text{(기간 합 플랫폼 이윤)}\label{c:profit}
\end{align}

\begin{itemize}[leftmargin=1.5em]
  \item 식 \eqref{c:bundle}: $y_b=1$이면 모든 품목의 요구수량을 정확히 채워야 하고,
        $y_b=0$이면 모든 배정이 0이다. 여러 공급자의 $x_{bji}$를 더하므로 분할 공급이 가능하다.
        이 한 식이 ``모든 품목이 확보되어야 성립''과 ``분할 공급 허용''을 함께 표현한다.
  \item 식 \eqref{c:cap}: 공급자별·품목별 용량을 넘지 않는다.
  \item 식 \eqref{c:rule}--\eqref{c:free}: 규칙 기반은 잉여를 마진의 고정 비율로 묶고,
        근시안과 가치 근사 정책은 자기 거래 마진 범위 안에서 자유롭게 고른다.
  \item 식 \eqref{c:profit}: 개별 주문이 아니라 기간 합에 대한 하나의 제약이다(교차 보조 허용).
\end{itemize}

\subsection{정책은 이 MILP의 특수한 경우}

\begin{table}[H]
\centering\small
\begin{tabularx}{\linewidth}{lYY}
\toprule
정책 & 잉여 & 유지 가치\\
\midrule
근시안 & 식 \eqref{c:free} (자유) & $c=0$ $\Rightarrow$ 잉여를 모두 플랫폼이 가져감\\
규칙 기반 & 식 \eqref{c:rule} (고정 비율) & $c=0$\\
선형 / MLP / GNN 근사 & 식 \eqref{c:free} (자유) & 학습한 $\hat V$로 계산한 $c_n$\\
\bottomrule
\end{tabularx}
\end{table}

근시안은 오늘 이윤만 보므로 주문자·공급자에게 잉여를 주지 않는다.
따라서 재참여 확률이 바닥값(0.5)에 머무는 ``미래를 전혀 고려하지 않는'' 참고용 하한이다.

\subsection{정산}

솔버 해의 배정 $x$는 정수 허용오차($10^{-5}$) 때문에 반올림한 뒤 쓰고,
잉여는 반올림한 배정으로 다시 계산한 마진에 맞춘다(규칙 기반은 고정 비율 그대로,
그 외는 $[0,\text{마진}]$으로 자름). 환경은 결정을 그대로 믿지 않고 묶음 성립, 용량,
정수성, 잉여 범위, 플랫폼 이윤 비음수를 다시 검사한다.

% ============================================================
\section{가치함수 근사 알고리즘}
% ============================================================

\subsection{전체 절차}

\begin{enumerate}[leftmargin=1.5em]
  \item \textbf{시장 만들기}: 시장 조건으로 주문자·공급자 프로필과 공급 구조를 생성한다.
  \item \textbf{경험 모으기}: 튜닝 seed에서 \emph{규칙 기반 정책}으로 시뮬레이션하며
        (상태, 이후 실제로 번 누적 이윤) 쌍을 저장한다.
  \item \textbf{시뮬레이션 기준치 만들기}: 저장된 상태 일부에서 같은 후속 정책(규칙 기반)과
        공통 난수로 여러 번 롤아웃해 평균 누적 이윤을 구한다. 예측 모델을 채점하는 기준이다.
  \item \textbf{모델 학습}: 선형 · MLP (· MLP+편중도 요약 · GNN)를 같은 데이터, 같은 검증 분할,
        같은 튜닝 예산(설정 5개)으로 학습한다.
  \item \textbf{유지 가치 변환}: 학습한 $\hat V$로 참여자별 유지 가치 $c_n$을 계산한다.
  \item \textbf{정책 평가}: 평가 seed에서 같은 MILP에 유지 가치만 바꿔 넣어 30반복 운영하고,
        누적 이윤·충족률·유지율을 짝지은 차이와 95\% 신뢰구간으로 비교한다.
\end{enumerate}

\subsection{입력 표현}

\textbf{시장 요약 지표} (선형·MLP 공통, 34개): 주문자 수, 공급자 수, 품목별 총 요구량,
품목별 총 공급가능량, 품목별 평균 제안가격, 품목별 평균 공급가격, 품목별 대체 공급자 수,
주문자·공급자 평균 재참여 확률 등. 선형과 MLP는 반드시 이 같은 지표만 쓴다.

\textbf{거래 연결 구조} (GNN): 주문자와 공급자를 노드로 두고, 주문자가 원하는 품목을
공급자가 하나 이상 공급할 수 있으면 연결선을 긋는다.
\begin{itemize}[leftmargin=1.5em]
  \item 주문자 노드: 품목별 요구수량, 품목별 제안가격, 재참여 확률
  \item 공급자 노드: 품목별 공급가능량, 품목별 가격, 재참여 확률
  \item 참여자 수가 기간마다 달라지므로 최대 참여자 수만큼 자리를 두고 활성 마스크로 처리
  \item 연결선 정보는 (1) 거래 가능 여부 $\to$ (2) 품목별 거래 가능 여부 $\to$ (3) 품목별 잠재 잉여로
        단계적으로 확장하고, 확장할 때마다 시장 요약 지표에도 같은 정보의 합계·평균을 추가한다.
\end{itemize}

\textbf{MLP + 편중도 요약}: 시장 요약 지표에 공급자별 품목 수의 최댓값·분산을 더한 비교 모델.
GNN이 이 모델보다 나아야 ``숫자 몇 개로는 요약되지 않는 구조''를 본다고 말할 수 있다.

\subsection{유지 가치 변환}

\begin{equation}
  c_n = \hat V(s) - \hat V(s\setminus n)
\end{equation}
참여자 $n$을 뺀 상태를 다시 예측해 그 차이를 유지 가치로 삼는다. 선형·MLP는 $n$을 뺀
시장 요약 지표로, GNN은 노드 $n$을 지운 그래프로 다시 예측한다. 예측이 튀는 것을 막기 위해
전체 평균 쪽으로 줄이고(shrink 0.5), 상한을 둔다(clip $0.2\,\hat V(s)$). 이 가드는
세 근사 방법에 똑같이 적용한다.

\subsection{해석: 한 단계 정책 개선}

경험과 기준치를 규칙 기반 정책으로 만들었으므로, 모델이 배우는 $\hat V(s)$는
``앞으로 \emph{규칙 기반 정책으로 운영한다면} 얼마나 벌까''이다.
이 가치를 MILP에 넣으면 ``오늘은 미래를 고려해 결정하고, 내일부터는 규칙 기반으로 간다''는
계산이 매 기간 반복된다. 즉 이 방법은 기준 정책을 한 단계 개선하는 근사 동적계획법(ADP)이다.
개선된 정책으로 다시 데이터를 모아 학습을 반복하면 더 멀리 보는 정책을 만들 수 있으나,
이번 예비 실험에서는 한 번만 수행한다.

% ============================================================
\section{실험 설계}
% ============================================================

\subsection{세 실험}

\textbf{실험 1: 선형 근사의 한계와 학습된 근사의 필요성}
\begin{itemize}[leftmargin=1.5em]
  \item 질문: 같은 시장 요약 지표를 줄 때, 선형 근사(기존 ADP)와 신경망 근사의
        예측 정확도와 장기 성과는 어떻게 다른가?
  \item 검증할 가설: 지표들 사이에 조합 효과가 있다면 신경망의 예측 오차가 더 작고 장기 성과가 높다.
  \item 반대 결과의 의미: 차이가 없으면 이 환경의 미래 가치는 선형으로 충분히 설명된다.
  \item 측정: 시뮬레이션 기준치 대비 예측 오차, 유지 가치 변환 오차, 누적 이윤·충족률·유지율.
        근시안을 항상 함께 표시하며, 신경망이 근시안을 유의하게 넘는지가 실질적인 질문이다.
\end{itemize}

\textbf{실험 2: 반응 함수에 대한 강건성} (이번 범위에서 보류)
\begin{itemize}[leftmargin=1.5em]
  \item 질문: 재참여 반응이 로지스틱이 아닐 때(예: 계단형) 성능은 어떻게 변하는가?
  \item 두 경우: 바뀐 반응에 맞춰 재학습 / 로지스틱으로 학습하고 다른 반응 환경에서 평가.
\end{itemize}

\textbf{실험 3: 구조 정보의 유효 조건}
\begin{itemize}[leftmargin=1.5em]
  \item 질문: 시장 요약 지표만 쓰는 MLP와 거래 연결 구조를 반영한 GNN의 성과는
        어떤 시장 조건에서 달라지는가?
  \item 검증할 가설: 공급 편중도가 높고, 대체 공급자가 적고, 주문당 품목 수가 많을수록
        구조 정보의 효과가 커진다.
  \item 대조 조건: \textbf{주문당 품목 수 1}. 공급자 구조는 그대로 두고 주문만 단품으로 바꾼 시장이다.
        묶음이 없으면 공급자가 떠날 때의 손해가 품목별로 따로 더해지므로, 시장 요약 지표만으로
        충분하고 두 방법의 차이는 통계적으로 구분되지 않아야 한다.
  \item 공급 편중도 0은 대조 조건이 아니라 \emph{기준 수준}이다. 모든 공급자가 같은 수의 품목을 맡아도
        어떤 품목끼리 같은 공급자에 묶이는지가 남아 있기 때문이다. 가설은
        ``편중도를 올릴수록 효과가 커진다''는 추세로 검증한다.
  \item 핵심 설계: $n_{\mathrm{alt}}=2$ 고정, 공급 편중도 $\{0,0.5,1\}\times$ 주문당 품목 수 $\{1,2,3\}$.
  \item 보조 설계: $n_{\mathrm{alt}}=3$ (평균 품목 수가 함께 바뀌어 효과가 섞임을 명시).
  \item 비교: MLP vs MLP+편중도 요약 vs GNN.
\end{itemize}

\subsection{비교 정책}

\begin{itemize}[leftmargin=1.5em]
  \item \textbf{근시안}: 유지 가치 0. 미래를 전혀 고려하지 않을 때의 참고용 하한.
  \item \textbf{규칙 기반}: 잉여를 사전에 정한 고정 비율로 배분. 주문자 몫 $\{0,0.1,0.2,0.3\}$
        $\times$ 공급자 몫 $\{0.1,0.2,0.3,0.4\}$ 격자에서 ``기준값 $(0.3,0.3)$보다 유의하게
        나을 때만 바꾼다''는 규칙으로 튜닝 seed에서 골랐고, 결과는 $(0.3,0.3)$이다.
  \item \textbf{선형 가치 근사}(기존 ADP), \textbf{MLP}, \textbf{MLP+편중도 요약}, \textbf{GNN}.
\end{itemize}

\subsection{공정성과 편향 방지 원칙}

\begin{enumerate}[leftmargin=1.5em]
  \item 비교군을 약하게 만들어 제안 방법이 좋아 보이게 하지 않는다.
  \item 모든 정책은 같은 MILP를 공유하고 유지 가치만 다르다(유지 가치 0 = 근시안, 테스트로 확인).
  \item 모든 근사 방법에 같은 입력 정보(해당 실험에서), 같은 학습 데이터·검증 분할, 같은 튜닝 예산을 준다.
  \item 튜닝·학습 seed와 평가 seed를 분리하고, 최종 비교는 평가 seed로 한 번만 수행한다.
  \item 가설은 검증 대상일 뿐 기대 결과가 아니다. 반대 결과도 동등하게 유효한 결과로 기록한다.
  \item 실험 설정(규모, 반복 수, 탐색 범위)은 결과를 보기 전에 고정하고,
        결과를 본 뒤의 변경은 모두 이유와 함께 결정 기록에 남긴다.
  \item 모든 표와 그림은 \texttt{result/runs/*.json} 원자료에서만 만든다.
\end{enumerate}

% ============================================================
\section{설계가 바뀌어 온 과정}
% ============================================================

실험 설계는 처음부터 지금 모습이 아니었다. 구현하면서 드러난 문제를 하나씩 바로잡으며
현재 형태가 되었다. 아래 표는 그 과정을 순서대로 정리한 것이다.

{\small
\begin{longtable}{p{0.5cm}p{4.4cm}p{4.6cm}p{4.9cm}}
\toprule
& 드러난 문제 & 결정 & 이유\\
\midrule
\endhead
1 & 공급자–품목 그래프에서 주문 쪽 정보가 합계로 뭉개져 ``함께 주문되는 품목''이 사라짐 &
주문자와 공급자를 모두 노드로 두는 거래 연결 구조로 변경 &
참여자가 곧 노드가 되어 유지 가치를 ``노드 하나 지우기''로 자연스럽게 계산 가능\\
2 & ``다품목 공급자 비율''을 올리면 공급자 수가 줄어 두 효과가 섞임 &
공급자 수·품목 수·대체 공급자 수를 모두 고정하고 \textbf{공급 편중도}만 바꿈 &
시장 요약 지표가 같고 구조만 다른 시장을 만들어야 초록의 주장을 직접 검증\\
3 & 편중도 0에서도 ``어떤 품목끼리 묶이는지''라는 구조 정보가 남음 &
대조 조건을 \textbf{주문당 품목 수 1}로 변경, $n_{\mathrm{alt}}=1$은 제외 &
단품 주문 시장은 구조 정보가 쓸모없으면서 시장 자체는 정상. $n_{\mathrm{alt}}=1$은 대체 공급자가 없는 극단적 시장\\
4 & $k$를 바꾸면 품목별 수요가 달라짐 &
주문자 수를 고정하고 품목당 수량을 $k_{\mathrm{ref}}/k$배로 조정, 수량 단위 3배 &
주문자 수(노드 수)를 조건 간 같게 유지, $k=3$의 반올림 문제 완화\\
5 & 큰 단품 주문을 공급자 한 명이 채우지 못함 &
한 품목의 \textbf{분할 공급 허용} &
허브 통합 배송 구조에 맞고, $k$별 주문 채우기 난이도 차이 제거\\
6 & 배송비·기회손실 처리 방식 &
주문당 고정비만 부과, 기회손실은 기록만 &
모델 단순화, 이중 계산 방지\\
7 & 주문 한 건이 품목 용량에 비해 너무 커 공급 부족 거절 19\% &
신규 주문자 수와 공급용량을 같은 비율로 확대 (주문 1건 $\le$ 용량의 1/4) &
공급 여유 비율은 유지하면서 ``덩어리'' 문제 해소\\
8 & 근시안은 잉여를 모두 가져가 참여자가 거의 모두 떠남 &
경험 수집·기준치의 후속 정책을 규칙 기반으로 변경, 근시안은 참고용 하한 &
근시안 기준으로 배우면 ``아무도 남지 않는 미래''만 학습하게 됨\\
9 & 공급자 점유율이 설계보다 낮아 공급이 얇음 &
용량을 공급자 점유율로 나눠 잡고, $T$를 25$\to$30으로 &
주문자 쪽과 같은 안정 상태 원칙, 유지 투자 회수 기간 확보\\
10 & 빈자리 즉시 충원은 비현실적 &
충원 확률 0.5$\to$0.25 &
새 공급자 확보에 시간이 걸리는 현실 반영, 공급자 유지 가치 성립\\
11 & 규칙 기반의 주문자 몫 0.3이 임의값 &
주문자 몫도 공급자 몫과 함께 격자 튜닝 &
비교군을 약하게 두지 않기 위해\\
12 & 잉여 0에서 재참여 0.1이면 거래 못 한 참여자의 90\%가 떠나 시장이 무너짐 &
재참여 곡선을 두 점($P(0)=0.5$, $P(\text{기준})=0.85$)으로 직접 정하고, 용량은 실측 점유율로 &
악순환(빈 품목$\to$거래 실패$\to$이탈$\to$더 많은 빈 품목) 제거\\
13 & 10반복의 신뢰구간이 몫 간 차이보다 큼 &
평가 반복 수 10$\to$30 &
검정력 확보. 이 과정에서 10반복에서 좋아 보이던 후보가 선택 효과였음을 발견\\
14 & 이윤 제약 단위와 잉여율 분모의 의미 &
기간 합 제약(교차 보조 허용) 유지, 공급자 분모는 내놓은 공급 전체임을 명시 &
손해를 감수하고 핵심 참여자를 붙잡는 것이 연구가 보려는 행동\\
\bottomrule
\end{longtable}
}

\keypoint{\textbf{원칙.} 환경 변경 중 일부(7--13)는 비교군 결과를 본 뒤에 이루어졌다.
이 변경들은 제안 방법을 유리하게 하려는 것이 아니라 ``참여자 유지의 가치'' 자체가 성립하는
시장을 만들기 위한 것이며, 모두 결정 기록에 변경 전 결과와 함께 남아 있다.
제안 방법(가치 근사)의 결과를 본 뒤에는 환경을 바꾸지 않았다.}

% ============================================================
\section{지금까지의 결과}
% ============================================================

모든 수치는 평가 seed 30반복(단, 표시한 곳 제외), $T=30$, 기본 시장 조건 기준이며,
$\pm$는 95\% 신뢰구간, ``짝지은 차이''는 같은 반복끼리 비교한 차이다.

\subsection{비교군 진단}

규칙 기반 $(0.3,0.3)$은 근시안보다 유의하게 낮았다(짝지은 차이 $-3978\pm1498$).
10반복에서 근시안보다 높아 보였던 작은 몫 후보 $(0.1,0.2)$, $(0.0,0.2)$는 새 반복(10--29)에서
각각 $-474$, $-411$로 돌아가, 표본이 적을 때의 선택 효과였음이 드러났다.

\subsection{실험 1: 모델 적합}

\begin{table}[H]
\centering\small
\caption{적합 결과 (튜닝 seed, 학습 630 / 검증 126)}
\begin{tabular}{lccc}
\toprule
모델 & 선택된 설정 & 격자 끝 여부 & 검증 $R^2$\\
\midrule
선형 & $\lambda=100$ & 예 (가장 강한 규제) & 0.347\\
MLP & weight decay $=0.1$ & 예 (가장 강한 규제) & 0.208\\
\bottomrule
\end{tabular}
\end{table}

두 모델 모두 가장 강한 규제를 골랐고, MLP는 그래도 선형보다 검증 오차가 컸다.
학습 표본 630개로는 신경망이 과적합한다는 신호다. 격자는 결과를 보기 전에 정했으므로 넓히지 않았다.

\subsection{실험 1: 예측 오차와 유지 가치 변환 오차}

\begin{table}[H]
\centering\small
\caption{예측 오차 (24개 상태, 시뮬레이션 기준치 평균 5747)}
\begin{tabular}{lcc}
\toprule
모델 & 편향 & 평균 절대 오차\\
\midrule
선형 & $-160\pm551$ & 1129\\
MLP & $-191\pm575$ & 1164\\
\bottomrule
\end{tabular}
\end{table}

절대 오차의 짝지은 차이(선형 $-$ MLP)는 $-35\pm333$으로 구분되지 않으며,
MLP가 더 정확한 상태는 24개 중 11개였다.

\begin{table}[H]
\centering\small
\caption{유지 가치 변환 오차 (12개 상태, 30 롤아웃)}
\begin{tabular}{llccc}
\toprule
모델 & 대상 & 평균 절대 차이 & 모델의 $c$ & 기준치의 $c$ (잡음 수준)\\
\midrule
선형 & 주문자 & 74 & 82 & 66 (49)\\
선형 & 공급자 & 594 & 652 & 913 (195)\\
MLP & 주문자 & 92 & 40 & 66 (49)\\
MLP & 공급자 & 732 & 696 & 913 (195)\\
\bottomrule
\end{tabular}
\end{table}

두 모델 모두 오차가 잡음 수준보다 크고, \textbf{공급자의 유지 가치를 과소평가}한다(913 대비 650--700).

\subsection{실험 1: 정책 비교 (평가 seed에서 한 번만 수행)}

\begin{table}[H]
\centering\small
\caption{정책별 성과}
\begin{tabular}{lcccc}
\toprule
정책 & 재참여율 (주문자/공급자) & 공급자 없는 품목 & 충족률 & 누적 이윤\\
\midrule
근시안 & 0.501 / 0.506 & 1.97 & 0.358 & $\mathbf{20204}\pm1847$\\
규칙 기반 $(0.3,0.3)$ & 0.676 / 0.770 & 0.86 & 0.516 & $16226\pm2349$\\
선형 유지 가치 & 0.535 / 0.712 & 1.22 & 0.467 & $19247\pm1895$\\
MLP 유지 가치 & 0.534 / 0.660 & 1.43 & 0.437 & $18437\pm1906$\\
\bottomrule
\end{tabular}
\end{table}

\begin{table}[H]
\centering\small
\caption{누적 이윤의 짝지은 차이}
\begin{tabular}{lcc}
\toprule
비교 & 평균 $\pm$ 95\% CI & 양수인 반복\\
\midrule
선형 $-$ 규칙 기반 & $+3021\pm1732$ (유의) & 24/30\\
MLP $-$ 근시안 & $-1768\pm704$ (유의) & 5/30\\
MLP $-$ 선형 & $-810\pm1012$ & 10/30\\
선형 $-$ 근시안 & $-958\pm1118$ & 9/30\\
\bottomrule
\end{tabular}
\end{table}

\subsection{해석}

\begin{enumerate}[leftmargin=1.5em]
  \item \textbf{가설은 지지되지 않았다.} 예측 오차는 구분되지 않고, 검증 $R^2$는 MLP가 낮으며,
        정책 성과도 낫지 않다. 다만 두 모델 모두 가장 강한 규제를 고른 점으로 보아,
        ``신경망이 필요 없다''기보다 ``이 데이터 양(630)으로는 신경망의 장점이 드러나지 않는다''에 가깝다.
  \item \textbf{가치 근사는 기준 정책을 개선했다.} 선형 유지 가치 정책은 규칙 기반을 유의하게 넘었다
        ($+3021$). 이는 한 단계 정책 개선이라는 방법의 이론적 역할과 일치한다.
  \item \textbf{근시안을 넘지 못한 이유는 출발점이다.} 모델이 배운 가치는 ``규칙 기반으로 운영할 때의 미래''인데,
        규칙 기반 자체가 근시안보다 약하다. 약한 기준 정책을 한 번 개선한 것만으로는 근시안에 닿지 못했다.
  \item \textbf{공급자 유지 가치의 과소평가}도 원인이 될 수 있다. 가치 근사 정책은 공급자 재참여를
        0.51에서 0.66--0.71로 올렸지만, 공급자 없는 품목을 규칙 기반만큼 줄이지는 못했다.
\end{enumerate}

% ============================================================
\section{현재 위치와 다음 단계}
% ============================================================

구현은 ``한 번에 하나씩, 매 단계가 항상 돌아가는 상태''를 유지하는 슬라이스 단위로 진행했다.
각 슬라이스는 테스트 통과와 MILP 최적해 확인 뒤 커밋하고 멈춘다.

\begin{table}[H]
\centering\small
\caption{슬라이스별 진행 상황}
\begin{tabularx}{\linewidth}{lYl}
\toprule
슬라이스 & 내용 & 상태\\
\midrule
0 준비 & 설정, 공통 난수, 참여자 프로필, 시장 조건(공급 편중도·대체 공급자 수·주문당 품목 수) & 완료\\
1 최소 경로 & 환경, 수락·배정 MILP(분할 공급), 시뮬레이션 실행기, 공급 규모 조정 & 완료\\
2 재참여·잉여배분 & 재참여 곡선, 구간선형 근사, 잉여배분 MILP, 근시안·규칙 기반, 환경 진단 & 완료\\
3 가치 근사 한 바퀴 & 시장 요약 지표, 시뮬레이션 기준치, 경험 수집, 선형 근사, 유지 가치 변환 & 완료\\
4 MLP & 신경망 근사, \textbf{실험 1} & \textbf{완료 (현재 위치)}\\
5 GNN & 거래 연결 구조, GNN, MLP+편중도 요약, \textbf{실험 3} & 다음\\
6 반응 함수 & 계단형 반응, \textbf{실험 2} & 보류\\
7 결과 정리 & 원자료에서 그림·요약표 생성 & 예정\\
\bottomrule
\end{tabularx}
\end{table}

\subsection{슬라이스 5 전에 정할 것}

\begin{itemize}[leftmargin=1.5em]
  \item \textbf{학습 데이터 확대}: MLP가 630개 표본에서 과적합했으므로, 파라미터가 더 많은 GNN은
        같은 문제를 겪을 가능성이 크다. 슬라이스 5 결과를 보기 전에 학습 데이터를 늘릴지 정해야 하며,
        늘린다면 GNN·MLP·MLP+편중도 요약에 똑같이 적용하고 실험 1의 630개 결과는 그대로 둔다.
  \item \textbf{실험 3의 평가 지표}: 정책 성과가 근시안을 넘지 못하더라도, 예측 오차와 유지 가치 변환 오차로
        MLP·MLP+편중도 요약·GNN의 차이를 비교할 수 있다.
  \item \textbf{범위}: 시간이 부족하면 공급 편중도 $\{0,1\}\times$ 주문당 품목 수 $\{1,2\}$의 축소 설계로 진행한다.
\end{itemize}

\subsection{이후 과제}

\begin{itemize}[leftmargin=1.5em]
  \item 반복 정책 개선: 개선된 정책으로 다시 데이터를 모아 재학습하여 근시안을 넘는지 확인
  \item 실험 2(반응 함수 강건성), 실험 3 보조 설계($n_{\mathrm{alt}}=3$)
  \item 이윤 상한(사후 최적) 계산, 규칙 기반 비교군의 이분탐색 튜닝
  \item 재참여 수준(0.5 외)에 대한 민감도 분석
\end{itemize}

% ============================================================
\appendix
\section{용어 정리}
% ============================================================

\begin{tabularx}{\linewidth}{lY}
\toprule
용어 & 뜻\\
\midrule
시장 조건 & 실험 3에서 바꾸는 세 값: 공급 편중도, 대체 공급자 수, 주문당 품목 수\\
공급 편중도 & 공급자들이 품목을 얼마나 고르게 나눠 갖는지 (0 = 고르게, 높을수록 소수에게 몰림)\\
시장 요약 지표 & 주문자·공급자 정보를 합계·평균·개수로 요약한 값 (선형·MLP의 입력)\\
거래 연결 구조 & 어느 주문자가 어느 공급자와 어떤 품목으로 거래할 수 있는지의 연결 관계 (GNN의 입력)\\
유지 가치 & 참여자 한 명을 붙잡아 두는 것의 미래 가치, $c_n=\hat V(s)-\hat V(s\setminus n)$\\
시뮬레이션 기준치 & 같은 상태에서 같은 후속 정책으로 여러 번 시뮬레이션해 얻은 평균 누적 이윤. 예측 모델의 오차를 재는 기준\\
재참여 확률 & 배분 잉여율에 따른 다음 기간 참여 확률 (로지스틱)\\
근시안 정책 & 유지 가치를 0으로 둔 정책 (참고용 하한)\\
규칙 기반 정책 & 잉여를 사전에 정한 고정 비율로 배분하는 정책\\
슬라이스 & 환경$\to$정책$\to$평가를 관통하는 얇은 구현 단위\\
\bottomrule
\end{tabularx}

\end{document}
